\subsection{Summary}
  \label{subsec:void_summary}

  In this section we have shown that the same effective saturation mechanism
  that accounts for galaxy rotation curves naturally extends to
  late-time, low-density cosmological environments.

  When the characteristic gravitational field falls below the
  empirically determined scale $a_\star$,
  the effective description operates in a bounded-response regime.
  In cosmic voids, this regime induces an environment-dependent
  bias in the locally inferred expansion rate,
  while leaving the global background dynamics unchanged.

  The mechanism therefore provides a structural link between
  galactic kinematics and the Hubble tension:
  both phenomena arise from the same transition scale,
  independently fixed by rotation curve data.

  The predicted amplitude of the effect is constrained to the
  percent-to-few-percent level,
  depends on independently measured void density contrasts,
  and decreases with increasing redshift due to environmental averaging.
  These features render the proposal observationally testable.

  In the next section, we examine the robustness of these results,
  discuss degeneracies with standard cosmological effects,
  and clarify the domain of validity of the effective description.
